\documentclass{computacion}
\titulo{Manual de Instalación --- Sistema CyAD}
\tipo{Proyecto Tecnológico}
\version{Manual de Instalación}
\trimestre{2026-Primavera}
\alumno{[TU NOMBRE]}
\matricula{[MATRÍCULA]}

\asesor{[Grado. Nombres Apellidos del Asesor]}
\categoria{[Profesor Titular/Asociado/Asistente]}
\departamento{[Departamento de Sistemas]}
\correo{[correo@azc.uam.mx]}

\coasesor{}
\categoriacoasesor{}
\departamentocoasesor{}
\correocoasesor{}

\begin{document}
\portada
\setcounter{page}{1}
\pagestyle{fancy}

\section{Introducción}
Este manual describe cómo instalar y ejecutar el \emph{Sistema para la gestión de formularios sobre los cursos y autoevaluaciones docentes} de la División de Ciencias y Artes para el Diseño (CyAD) en un entorno de desarrollo local. La instalación en producción sigue los mismos pasos, cambiando únicamente el archivo \texttt{.env} y el modo de servido (\emph{gunicorn} + \emph{Nginx}).

\section{Prerrequisitos}

\begin{longtable}{p{0.30\textwidth} p{0.20\textwidth} p{0.42\textwidth}}
    \toprule
    \textbf{Software} & \textbf{Versión mínima} & \textbf{Uso} \\
    \midrule
    \endhead
    Python & 3.11+ (probado con 3.13) & Ejecución del backend Django. \\
    \midrule
    Node.js & 20+ (probado con 22 LTS) & Ejecución del frontend Vite + React. \\
    \midrule
    PostgreSQL & 14+ (probado con 17) & Base de datos relacional. \\
    \midrule
    Git & 2.40+ & Clonado del repositorio. \\
    \midrule
    \emph{npm} (incluido con Node) & 10+ & Gestor de paquetes JS. \\
    \bottomrule
\end{longtable}

En Windows se recomienda utilizar \emph{Git Bash} o \emph{PowerShell}. En Linux/macOS los comandos son idénticos, cambiando la ruta del entorno virtual (\texttt{bin/} en lugar de \texttt{Scripts/}).

\textbf{Nota importante para Windows:} el instalador oficial de PostgreSQL no agrega \texttt{psql} al \texttt{PATH} del sistema. Si al escribir \texttt{psql --version} obtiene ``command not found'', agregue manualmente la carpeta \texttt{C:\textbackslash Program Files\textbackslash PostgreSQL\textbackslash 17\textbackslash bin} (ajuste al número de versión que tenga) a la variable de entorno \texttt{Path} de Windows, o invoque \texttt{psql} por ruta completa.

\section{Clonado del repositorio}

\begin{quote}
\ttfamily
git clone \textless URL\_del\_repositorio\textgreater\ proyecto\_terminal\_cyad \\
cd proyecto\_terminal\_cyad
\end{quote}

\section{Configuración del backend}

\subsection{Entorno virtual}

\begin{quote}
\ttfamily
cd backend \\
python -m venv .venv \\
.venv\textbackslash Scripts\textbackslash activate\ \ \# Windows \\
source .venv/bin/activate\ \ \# Linux/macOS \\
pip install -r requirements.txt
\end{quote}

\subsection{Base de datos}

El proyecto utiliza por defecto el usuario \texttt{postgres} y una base de datos llamada \texttt{proyecto\_terminal\_cyad} (ver \texttt{backend/.env.example}). La forma recomendada es dejar que el script auxiliar cree la base de datos automáticamente después de configurar el archivo \texttt{.env} (siguiente sección):

\begin{quote}
\ttfamily
python scripts/create\_db.py
\end{quote}

Si prefiere crearla manualmente con \texttt{psql}:

\begin{quote}
\ttfamily
psql -U postgres \\
CREATE DATABASE proyecto\_terminal\_cyad;
\end{quote}

También puede definir un usuario dedicado con permisos limitados (recomendado en despliegues compartidos) y ajustar \texttt{DB\_USER}/\texttt{DB\_PASSWORD} en el archivo \texttt{.env} en consecuencia:

\begin{quote}
\ttfamily
CREATE USER cyad\_user WITH PASSWORD 'cyad\_pass'; \\
GRANT ALL PRIVILEGES ON DATABASE proyecto\_terminal\_cyad TO cyad\_user;
\end{quote}

\subsection{Archivo \texttt{.env}}

\textbf{Ubicación:} el archivo \texttt{.env} que Django lee debe estar en la carpeta \texttt{backend/} (donde vive \texttt{manage.py}), \emph{no} en la raíz del repositorio. El repositorio incluye \texttt{backend/.env.example} como plantilla; cópielo y ajuste los valores:

\begin{quote}
\ttfamily
cd backend \\
cp .env.example .env
\end{quote}

Contenido esperado (valores por defecto de la plantilla; ajuste \texttt{DB\_PASSWORD} y \texttt{SECRET\_KEY} a su entorno):

\begin{Verbatim}[fontsize=\small, frame=single]
DEBUG=True
SECRET_KEY=cambia-esta-clave-por-una-aleatoria-larga
ALLOWED_HOSTS=localhost,127.0.0.1
DB_NAME=proyecto_terminal_cyad
DB_USER=postgres
DB_PASSWORD=tu_password_de_postgres
DB_HOST=localhost
DB_PORT=5432
JWT_ACCESS_MINUTES=60
JWT_REFRESH_DAYS=7
CORS_ALLOWED_ORIGINS=http://localhost:5173,http://127.0.0.1:5173
\end{Verbatim}

\textbf{Nota:} existe también un archivo \texttt{.env.example} en la raíz del repositorio; \emph{no} es el que lee Django y su esquema de variables JWT (\texttt{ACCESS\_TOKEN\_LIFETIME}, \texttt{REFRESH\_TOKEN\_LIFETIME}) no coincide con lo que espera \texttt{backend/config/settings.py}. Utilice siempre \texttt{backend/.env.example}.

Para producción se debe establecer \texttt{DEBUG=False}, definir \texttt{ALLOWED\_HOSTS} con el dominio real y generar una \texttt{SECRET\_KEY} con \texttt{python -c "from django.core.management.utils import get\_random\_secret\_key; print(get\_random\_secret\_key())"}.

\subsection{Migraciones}

\begin{quote}
\ttfamily
python manage.py migrate
\end{quote}

\subsection{Datos iniciales}

La forma más rápida de dejar el sistema listo para usar es ejecutar el \emph{seed} de demostración, que es idempotente y crea catálogos, UEAs, usuarios de prueba, periodo activo y el formulario de autoevaluación en un solo comando:

\begin{quote}
\ttfamily
python scripts/seed\_demo.py
\end{quote}

De forma alternativa, se pueden cargar los catálogos base y los usuarios mínimos por separado (útil cuando ya existen datos):

\begin{quote}
\ttfamily
python manage.py loaddata catalogos/fixtures/departamentos.json \\
python manage.py loaddata catalogos/fixtures/licenciaturas.json \\
python manage.py loaddata catalogos/fixtures/posgrados.json \\
python manage.py loaddata catalogos/fixtures/areas.json \\
python manage.py loaddata catalogos/fixtures/periodos.json \\
python scripts/seed\_users.py
\end{quote}

Al finalizar, existirán las siguientes cuentas:

\begin{longtable}{p{0.30\textwidth} p{0.35\textwidth} p{0.30\textwidth}}
    \toprule
    \textbf{Rol} & \textbf{Correo} & \textbf{Contraseña} \\
    \midrule
    \endhead
    Administrador & admin@cyad.uam.mx & Admin1234! \\
    \midrule
    Profesor & profesor@cyad.uam.mx & Profesor1234! \\
    \bottomrule
\end{longtable}

Estas contraseñas deben cambiarse antes del despliegue a producción.

\subsection{Servidor de desarrollo del backend}

\begin{quote}
\ttfamily
python manage.py runserver 8000
\end{quote}

El backend queda disponible en \texttt{http://localhost:8000}. La documentación interactiva de la API está en \texttt{http://localhost:8000/api/docs/}.

\section{Configuración del frontend}

En una terminal aparte:

\begin{quote}
\ttfamily
cd frontend \\
npm install \\
npm run dev
\end{quote}

El frontend queda disponible en \texttt{http://localhost:5173}. Vite realiza \emph{hot reload} automático al editar archivos fuente.

Si el backend no está en \texttt{localhost:8000}, cree \texttt{frontend/.env.local} con:

\begin{Verbatim}[fontsize=\small, frame=single]
VITE_API_URL=http://tu-servidor:8000/api/v1
\end{Verbatim}

\textbf{Nota:} el nombre correcto de la variable es \texttt{VITE\_API\_URL} (verificable en \texttt{frontend/src/api/client.js}).

\section{Verificación}

\begin{enumerate}
    \item Abra \texttt{http://localhost:5173/} — debe mostrarse la vista pública ``Explorar''.
    \item Inicie sesión con \texttt{admin@cyad.uam.mx}. Debe redirigir al \emph{dashboard} de administrador.
    \item Cierre sesión, inicie con \texttt{profesor@cyad.uam.mx}. Debe redirigir al \emph{dashboard} de profesor.
    \item Ejecute la suite de pruebas: \texttt{pytest} desde \texttt{backend/}. Todas deben pasar.
\end{enumerate}

\section{Comandos administrativos útiles}

\begin{longtable}{p{0.50\textwidth} p{0.42\textwidth}}
    \toprule
    \textbf{Comando} & \textbf{Uso} \\
    \midrule
    \endhead
    \texttt{python manage.py cargar\_catalogos\_csv <archivo>} & Importar catálogos desde CSV. \\
    \midrule
    \texttt{python manage.py createsuperuser} & Crear cuenta administradora adicional. \\
    \midrule
    \texttt{python manage.py dumpdata <app>} & Exportar datos a JSON. \\
    \midrule
    \texttt{pytest --cov=. --cov-report=html} & Ejecutar pruebas con cobertura. \\
    \midrule
    \texttt{npm run build} (frontend) & Compilar producción estática (dist/). \\
    \bottomrule
\end{longtable}

\section{Solución de problemas}

\begin{itemize}
    \item \textbf{Error \texttt{FATAL: password authentication failed}}: verifique \texttt{DB\_USER} y \texttt{DB\_PASSWORD} en \texttt{.env} y que el usuario tenga permisos sobre la base de datos.
    \item \textbf{Error CORS al iniciar sesión desde el frontend}: verifique que \texttt{CORS\_ALLOWED\_ORIGINS} incluya \texttt{http://localhost:5173}.
    \item \textbf{\texttt{429 Too Many Requests} al iniciar sesión}: se activó el \emph{rate limiting} (5 intentos por minuto). Espere 60 segundos y reintente.
    \item \textbf{Puerto 5173 ocupado}: Vite usará automáticamente el siguiente puerto libre; verifique la URL impresa al inicio.
    \item \textbf{\texttt{seed\_demo.py} falla con ``La suma de pesos de las secciones debe ser 100\,\%'' }: en algunas versiones del script las secciones se crean sin \texttt{peso}, lo que hace fallar la publicación del formulario. Como \emph{workaround} temporal, edite \texttt{backend/scripts/seed\_demo.py} y añada \texttt{peso=Decimal("50")} a la creación de cada una de las dos \texttt{Seccion.objects.create(\ldots)}, luego ejecute \texttt{python manage.py flush --noinput} y vuelva a lanzar el seed. Ver \texttt{hallazgos\_despliegue.pdf} para detalles.
    \item \textbf{\texttt{psql: command not found} en Windows}: agregue \texttt{C:\textbackslash Program Files\textbackslash PostgreSQL\textbackslash<versión>\textbackslash bin} al \texttt{PATH}, o llame a \texttt{psql} por ruta completa. Alternativamente, use \emph{pgAdmin} o el script \texttt{scripts/create\_db.py}, que utiliza el driver Python y no requiere \texttt{psql} en el PATH.
\end{itemize}

\end{document}
